\documentclass[aspectratio=169,10pt]{beamer}
\usepackage{fontspec}
\usetheme{Antibes}
\setbeamertemplate{navigation symbols}{}

% Page number bottom-right
\setbeamertemplate{footline}{
  \leavevmode%
  \hbox{%
    \begin{beamercolorbox}[wd=\paperwidth,ht=2.5ex,dp=1.2ex,right]{author in head/foot}%
      \usebeamerfont{author in head/foot}
      \insertframenumber{} / \inserttotalframenumber
      \hspace*{2ex}
    \end{beamercolorbox}%
  }%
}

% Custom headline: project name + section
\newcommand{\projectname}{APKGC: Noise-enhanced MMKGC with Attention Penalty (AAAI 2025)}
\setbeamertemplate{headline}{
  \leavevmode%
  \hbox{%
    \begin{beamercolorbox}[wd=\paperwidth,ht=2.5ex,dp=1.125ex,left]{section in head/foot}%
      \usebeamerfont{section in head/foot}\hspace*{2ex}\projectname%
    \end{beamercolorbox}}%
  \vskip0pt%
  \leavevmode%
  \hbox{%
    \begin{beamercolorbox}[wd=\paperwidth,ht=2.5ex,dp=1.125ex,left]{subsection in head/foot}%
      \usebeamerfont{subsection in head/foot}\hspace*{2ex}\insertsectionhead%
    \end{beamercolorbox}}%
  \vskip0pt%
}

% =====================
% Packages
% =====================
\usepackage{amsmath,amsfonts,amssymb}
\usepackage{graphicx}
\usepackage{booktabs}
\usepackage{tikz}
\usetikzlibrary{arrows.meta,positioning,fit,calc,shapes.geometric,backgrounds}
\usepackage{multirow}
\usepackage{colortbl}
\usepackage{xcolor}

% =====================
% Colors
% =====================
\definecolor{hcmusblue}{HTML}{1B3A5C}
\definecolor{contextblue}{HTML}{2E86C1}
\definecolor{methodteal}{HTML}{16A085}
\definecolor{objectivepurple}{HTML}{7D3C98}
\definecolor{riskred}{HTML}{C0392B}
\definecolor{lightbg}{HTML}{F5F8FC}

\definecolor{GoodGreen}{HTML}{1A6B3C}
\definecolor{GoodBg}{HTML}{D5EDE0}
\definecolor{BadRed}{HTML}{A82A1E}
\definecolor{BadBg}{HTML}{FAE1DE}
\definecolor{BlueBg}{HTML}{DBE9F8}
\definecolor{MutedTxt}{HTML}{555555}
\definecolor{OrangeAc}{HTML}{C05A00}

% Mã màu ba mô thái, dùng nhất quán toàn bài
\definecolor{structblue}{HTML}{4A7DC0}
\definecolor{visorange}{HTML}{C05A00}
\definecolor{textgreen}{HTML}{1A6B3C}

\setbeamercolor{structure}{fg=hcmusblue}
\setbeamercolor{block title}{fg=white,bg=contextblue}
\setbeamercolor{block body}{bg=contextblue!7}
\setbeamercolor{alerted text}{fg=riskred}

% =====================
% Semantic block environments
% =====================
\newenvironment{contextblock}[1]{%
  \setbeamercolor{block title}{fg=white,bg=contextblue}%
  \setbeamercolor{block body}{bg=contextblue!7}%
  \begin{block}{#1}}{\end{block}}

\newenvironment{methodblock}[1]{%
  \setbeamercolor{block title}{fg=white,bg=methodteal}%
  \setbeamercolor{block body}{bg=methodteal!7}%
  \begin{block}{#1}}{\end{block}}

\newenvironment{objectiveblock}[1]{%
  \setbeamercolor{block title}{fg=white,bg=objectivepurple}%
  \setbeamercolor{block body}{bg=objectivepurple!7}%
  \begin{block}{#1}}{\end{block}}

\newenvironment{riskblock}[1]{%
  \setbeamercolor{block title}{fg=white,bg=riskred}%
  \setbeamercolor{block body}{bg=riskred!7}%
  \begin{block}{#1}}{\end{block}}

% =====================
% Helper commands
% =====================
\newcommand{\emphkey}[1]{\textcolor{contextblue}{\textbf{#1}}}
\newcommand{\emphwarn}[1]{\textcolor{riskred}{\textbf{#1}}}
\newcommand{\emphok}[1]{\textcolor{methodteal}{\textbf{#1}}}
\newcommand{\emphobj}[1]{\textcolor{objectivepurple}{\textbf{#1}}}
\newcommand{\muted}[1]{\textcolor{MutedTxt}{#1}}

\newcommand{\es}{e_{\mathrm{s}}}
\newcommand{\ev}{e_{\mathrm{v}}}
\newcommand{\et}{e_{\mathrm{t}}}

% =====================
% Title
% =====================
\title[APKGC]{APKGC: Noise-enhanced Multi-Modal Knowledge Graph\\Completion with Attention Penalty}
\subtitle{Seminar GraphML --- Học phần HP3}
\author{Nhóm seminar}
\institute{ĐHQG-HCM --- Trường Đại học Khoa học Tự nhiên\\Khoa Công nghệ Thông tin}
\date{\today}
\setbeamerfont{title}{series=\bfseries}
\setbeamerfont{frametitle}{size=\normalsize,series=\bfseries}
\setbeamerfont{block title}{size=\footnotesize,series=\bfseries}
\setbeamerfont{block body}{size=\footnotesize}
\setlength{\medskipamount}{3pt}
\setlength{\bigskipamount}{6pt}

% Section dividers
\AtBeginSection[]{
  \begin{frame}
    \vfill\centering
    \begin{beamercolorbox}[sep=8pt,center,shadow=true,rounded=true]{title}
      \usebeamerfont{title}\insertsectionhead\par
    \end{beamercolorbox}
    \vfill
  \end{frame}
}

\begin{document}
\footnotesize

% -- TITLE ---------------------------------------------------------------------
\begin{frame}
  \titlepage
  \begin{center}
    \muted{\scriptsize AAAI 2025 (Oral), tr.\ 15005--15013 \quad$\cdot$\quad
           \texttt{github.com/HubuKG/APKGC}}
  \end{center}
\end{frame}

% -- OUTLINE -------------------------------------------------------------------
\begin{frame}{Nội dung}
  \tableofcontents
\end{frame}

% ==============================================================================
\section{Kiến thức nền}
% ==============================================================================

\begin{frame}{Đồ thị tri thức và bài toán hoàn thiện}
\begin{columns}[c]
\begin{column}{0.46\textwidth}
  \centering
  \begin{tikzpicture}[scale=0.72, transform shape,
    ent/.style={circle, draw=structblue, fill=BlueBg, minimum size=1.35cm,
                font=\tiny, align=center, line width=0.7pt},
    unk/.style={circle, draw=BadRed, fill=BadBg, dashed, minimum size=1.35cm,
                font=\large\bfseries, text=BadRed, line width=0.8pt},
    rel/.style={-{Stealth[length=2mm]}, structblue, line width=0.6pt},
    urel/.style={-{Stealth[length=2mm]}, BadRed, dashed, line width=0.6pt},
    lb/.style={font=\tiny, text=MutedTxt, midway, sloped, above}]
    \node[ent] (dh) at (0,2.0)   {Dennis\\Hopper};
    \node[ent] (bv) at (0,-1.4)  {Blue\\Velvet};
    \node[ent] (aw) at (3.2,0.3) {NSFC\\Award};
    \node[ent] (yr) at (6.4,2.0) {1987};
    \node[unk] (qq) at (6.4,-1.4){?};
    \draw[rel]  (dh) -- node[lb]{award}     (aw);
    \draw[rel]  (bv) -- node[lb]{nominated} (aw);
    \draw[rel]  (aw) -- node[lb]{year}      (yr);
    \draw[urel] (aw) -- node[lb,text=BadRed]{winner} (qq);
    \node[font=\tiny, text=BadRed] at (6.4,-2.5) {cần dự đoán};
  \end{tikzpicture}
\end{column}
\begin{column}{0.5\textwidth}
  \begin{contextblock}{Định nghĩa}
    Đồ thị tri thức $\mathcal{G}=(\mathcal{E},\mathcal{R},\mathcal{T})$.
    Mỗi sự kiện là một bộ ba $(h,r,t)$.
  \end{contextblock}

  \begin{riskblock}{Vấn đề gốc: tính không đầy đủ}
    KG do người xây \emphwarn{luôn thiếu}. DB15K có 12.842 thực thể nhưng chỉ
    79.222 bộ ba --- rất thưa so với không gian tổ hợp.
  \end{riskblock}

  \begin{objectiveblock}{Nhiệm vụ KGC}
    Cho $(h,r,?)$, xếp hạng toàn bộ $\mathcal{E}$ sao cho đáp án đúng lên đầu.
    Đây là bài toán \emphobj{xếp hạng}, không phải phân loại --- điều này quyết
    định cách chọn độ đo.
  \end{objectiveblock}
\end{column}
\end{columns}
\end{frame}

\begin{frame}{Nền tảng: từ TransE đến RotatE}
\begin{columns}[T]
\begin{column}{0.48\textwidth}
  \begin{contextblock}{TransE (NeurIPS'13) --- phép tịnh tiến}
    $h+r \approx t$, \quad $S=-\lVert h+r-t\rVert$
    \\[3pt]
    Quan hệ là một \emphkey{vector dịch chuyển}.
    \\[3pt]
    \emphwarn{Hạn chế:} không biểu diễn được quan hệ đối xứng --- nếu $r$ đối
    xứng thì buộc $r\approx 0$.
  \end{contextblock}
\end{column}
\begin{column}{0.48\textwidth}
  \begin{methodblock}{RotatE (ICLR'19) --- phép quay phức}
    $h\circ r \approx t$, \quad $\lvert r_i\rvert = 1$
    \\[3pt]
    Quan hệ là phép quay góc $\theta$ trong không gian phức.
    \\[3pt]
    \emphok{Xử lý được ba loại quan hệ khó:} đối xứng (quay $180^\circ$),
    nghịch đảo (quay ngược), hợp thành (cộng góc).
  \end{methodblock}
\end{column}
\end{columns}

\vspace{4pt}
\begin{objectiveblock}{Vì sao phải đi tiếp sang đa mô thái}
  Cả hai chỉ khai thác \emphobj{cấu trúc đồ thị}. Với thực thể xuất hiện ít
  (long-tail), cấu trúc gần như trống nên mô hình không đủ tín hiệu. Trong khi
  đó thực thể ngoài đời còn mang \emphobj{ảnh} và \emphobj{mô tả văn bản}.
  \textbf{APKGC kế thừa RotatE làm hàm chấm điểm.}
\end{objectiveblock}
\end{frame}

\begin{frame}{MMKG: mỗi thực thể mang ba mô thái}
\begin{columns}[T]
\begin{column}{0.55\textwidth}
  {\scriptsize
  \begin{tabular}{@{}llll@{}}
    \toprule
    Mô thái & Ký hiệu & Trích bằng & Trạng thái\\
    \midrule
    \textcolor{structblue}{Cấu trúc} & $\es$ & \texttt{nn.Embedding} & Huấn luyện\\
    \textcolor{visorange}{Thị giác}  & $\ev$ & BEiT, 4096-d & \emphwarn{Đóng băng}\\
    \textcolor{textgreen}{Văn bản}   & $\et$ & BERT, 768-d  & \emphwarn{Đóng băng}\\
    \bottomrule
  \end{tabular}}

  \vspace{6pt}
  \begin{methodblock}{Hệ quả của việc đóng băng}
    APKGC chỉ học \emphok{ma trận chiếu và cách hợp nhất}, không tinh chỉnh
    encoder. Đây là lý do mô hình chạy được trên một GPU đơn.
  \end{methodblock}
\end{column}
\begin{column}{0.42\textwidth}
  \begin{riskblock}{Thực tế: mô thái bị thiếu}
    Trong MKG-W nhiều thực thể \emphwarn{không có ảnh hoặc mô tả}. Repo lưu danh
    sách trong \texttt{missing.json}.
    \\[3pt]
    APKGC lấp chỗ trống bằng vector sinh từ $\mathcal{N}(\mu,\sigma)$ của các
    thực thể còn đủ mô thái --- hé lộ tinh thần xuyên suốt bài báo:
    \emphwarn{nhiễu không phải kẻ thù}.
  \end{riskblock}
\end{column}
\end{columns}
\end{frame}

% ==============================================================================
\section{Hai vấn đề bài báo đặt ra}
% ==============================================================================

\begin{frame}{Vấn đề 1 --- mô hình thiếu bền vững trước nhiễu}
\begin{columns}[T]
\begin{column}{0.48\textwidth}
  \begin{riskblock}{Cách làm trước đây}
    Xem nhiễu là yếu tố gây nhiễu loạn cần loại bỏ: lọc ảnh không khớp, loại mô
    tả sai lệch, làm sạch dữ liệu trước khi huấn luyện.
  \end{riskblock}

  \begin{riskblock}{Ba điểm hụt}
    \begin{itemize}\itemsep2pt
      \item Dữ liệu thực tế \emphwarn{luôn bẩn} --- không tồn tại tập triển khai đã sạch.
      \item Mô hình huấn luyện trên dữ liệu sạch \emphwarn{chưa từng gặp nhiễu}, dễ gãy khi vận hành.
      \item Lọc mạnh làm mất thông tin của nhóm thực thể hiếm --- đúng nhóm cần hỗ trợ nhất.
    \end{itemize}
  \end{riskblock}
\end{column}
\begin{column}{0.48\textwidth}
  \begin{methodblock}{Lật ngược của APKGC}
    Thay vì chống nhiễu, \emphok{chủ động tiêm nhiễu} trong quá trình huấn luyện.
  \end{methodblock}

  \begin{contextblock}{Liên hệ với kiến thức đã biết}
    Cùng nguyên lý với \emphkey{data augmentation} trong thị giác máy tính và
    với \emphkey{dropout}: buộc mô hình học đặc trưng bền vững thay vì ghi nhớ
    máy móc.
    \\[3pt]
    Dropout tắt ngẫu nhiên nơ-ron để mạng không phụ thuộc một đường dẫn duy
    nhất. APKGC làm điều tương tự ở \emphkey{mức mô thái}.
  \end{contextblock}
\end{column}
\end{columns}
\end{frame}

\begin{frame}{Vấn đề 2 --- attention over-trust}
\begin{columns}[c]
\begin{column}{0.44\textwidth}
  \centering
  \begin{tikzpicture}[scale=0.85, transform shape,
    lb/.style={font=\tiny, text=MutedTxt}]
    \node[font=\scriptsize\bfseries, text=BadRed, anchor=west] at (0,3.5)
         {Trước: chú ý dồn vào dấu chấm};
    \fill[gray!12] (0.00,1.85) rectangle (4.60,3.25);
    \fill[structblue!45] (0.05,1.9) rectangle ++(0.8,0.60);
    \fill[structblue!38] (1.00,1.9) rectangle ++(0.8,0.48);
    \fill[structblue!32] (1.95,1.9) rectangle ++(0.8,0.38);
    \fill[structblue!26] (2.90,1.9) rectangle ++(0.8,0.30);
    \fill[BadRed!85]     (3.85,1.9) rectangle ++(0.7,1.30);
    \node[lb] at (0.45,1.72) {Umbria};
    \node[lb] at (1.40,1.72) {is};
    \node[lb] at (2.35,1.72) {only};
    \node[lb] at (3.30,1.72) {region};
    \node[lb, text=BadRed] at (4.20,1.72) {.};
    \node[font=\tiny\bfseries, text=BadRed] at (4.20,3.38) {0.66};

    \node[font=\scriptsize\bfseries, text=GoodGreen, anchor=west] at (0,1.15)
         {Sau: phạt rồi tái phân bổ};
    \fill[gray!12] (0.00,-0.55) rectangle (4.60,0.85);
    \fill[GoodGreen!70] (0.05,-0.5) rectangle ++(0.8,0.85);
    \fill[GoodGreen!55] (1.00,-0.5) rectangle ++(0.8,0.62);
    \fill[GoodGreen!62] (1.95,-0.5) rectangle ++(0.8,0.72);
    \fill[GoodGreen!66] (2.90,-0.5) rectangle ++(0.8,0.78);
    \fill[GoodGreen!25] (3.85,-0.5) rectangle ++(0.7,0.26);
    \node[lb] at (0.45,-0.7) {Umbria};
    \node[lb] at (1.40,-0.7) {is};
    \node[lb] at (2.35,-0.7) {only};
    \node[lb] at (3.30,-0.7) {region};
    \node[lb] at (4.20,-0.7) {.};
    \node[font=\tiny\bfseries, text=GoodGreen] at (4.20,-0.10) {0.13};
  \end{tikzpicture}
\end{column}
\begin{column}{0.52\textwidth}
  \begin{riskblock}{Hiện tượng}
    Mô hình ngôn ngữ tiền huấn luyện dồn trọng số chú ý vào token
    \emphwarn{không mang ngữ nghĩa} --- dấu chấm, dấu phẩy, token đặc biệt.
    Trong ví dụ, dấu chấm nhận 0.66 còn danh từ riêng bị bỏ qua.
  \end{riskblock}

  \begin{contextblock}{Vì sao xảy ra}
    Được ghi nhận độc lập trong nghiên cứu về ảo giác đa mô thái
    (Huang et al., 2023). Token phổ biến xuất hiện ở mọi câu nên trở thành điểm
    trú ẩn an toàn khi mô hình không chắc chắn --- tương tự khái niệm
    \emphkey{attention sink}.
  \end{contextblock}

  \begin{methodblock}{Hướng chữa của APKGC}
    \emphok{Phát hiện và phạt} token bị tin tưởng quá mức, rồi \emphok{trả lại}
    phần chú ý đó cho token lân cận có ngữ nghĩa.
  \end{methodblock}
\end{column}
\end{columns}
\end{frame}

% ==============================================================================
\section{Kiến trúc APKGC}
% ==============================================================================

\begin{frame}{Toàn cảnh: ba khối nối tiếp}
\centering
\begin{tikzpicture}[scale=0.88, transform shape,
  blk/.style={draw=gray!50, fill=white, rounded corners=2pt, inner sep=4pt,
              text width=3.0cm, align=center, font=\tiny},
  ar/.style={-{Stealth[length=2mm]}, gray!70, line width=0.7pt}]

  \node[blk] (a1) at (0,3.0) {$\es$ \quad $\ev$ \quad $\et$};
  \node[blk, draw=visorange, fill=OrangeAc!8] (a2) at (0,2.0)
       {\textbf{Adaptive noise sampling}\\ $\rho>\theta_n$ ? trộn nhiễu : giữ nguyên};
  \node[blk] (a3) at (0,0.85)
       {$e \leftarrow (1{-}\text{mask})\,e + \text{mask}\cdot\text{noise}$};
  \begin{scope}[on background layer]
    \node[draw=visorange, rounded corners=4pt, inner sep=8pt,
          fill=OrangeAc!4, fit=(a1)(a3)] (ga) {};
  \end{scope}
  \node[font=\scriptsize\bfseries, text=visorange] at (0,4.15)
       {(a) Modality Noise Sampling};

  \node[blk] (b1) at (5.2,3.0) {Self-attention: $QK^{\!\top}\!/\sqrt{d}$};
  \node[blk, draw=BadRed, fill=BadBg] (b2) at (5.2,1.95)
       {\textbf{\textcolor{BadRed}{Attention Penalty}}\\
        tích luỹ hậu tố, phát hiện\\ và phạt token vượt ngưỡng $\phi$};
  \node[blk] (b3) at (5.2,0.85) {tái phân bổ cho cửa sổ $k$};
  \begin{scope}[on background layer]
    \node[draw=BadRed, rounded corners=4pt, inner sep=8pt,
          fill=BadBg!35, fit=(b1)(b3)] (gb) {};
  \end{scope}
  \node[font=\scriptsize\bfseries, text=BadRed] at (5.2,4.15)
       {(b) Entity Feature Extraction};

  \node[blk] (c1) at (10.4,3.0) {Modality interaction: $\theta_{ab}$};
  \node[blk] (c2) at (10.4,2.0) {hợp nhất ba mô thái};
  \node[blk, draw=GoodGreen, fill=GoodBg] (c3) at (10.4,0.85)
       {\textbf{\textcolor{GoodGreen}{RotatE score}}\\
        $S=\gamma-\lVert h\circ r-t\rVert$};
  \begin{scope}[on background layer]
    \node[draw=GoodGreen, rounded corners=4pt, inner sep=8pt,
          fill=GoodBg!35, fit=(c1)(c3)] (gc) {};
  \end{scope}
  \node[font=\scriptsize\bfseries, text=GoodGreen] at (10.4,4.15)
       {(c) Triple Score \& Prediction};

  \draw[ar] (ga.east) -- node[above, font=\tiny, text=MutedTxt]{token} (gb.west);
  \draw[ar] (gb.east) -- node[above, font=\tiny, text=MutedTxt]{$H$} (gc.west);
\end{tikzpicture}

\vspace{6pt}
\muted{\scriptsize Hai đóng góp chính nằm ở khối (a) và (b); khối (c) kế thừa RotatE.}
\end{frame}

\begin{frame}{Đóng góp 1 --- Adaptive Noise Sampling}
\begin{columns}[T]
\begin{column}{0.5\textwidth}
  \begin{methodblock}{Sinh nhiễu theo phân phối dữ liệu}
    Gaussian (dùng cho DB15K):
    \[ \mathbf{N}_g=\theta_r(\varphi_m\odot g+\mu_m),\quad g\sim K(0,Z) \]
    Average (dùng cho MKG-W):
    \[ \mathbf{N}_a=\theta_r[\beta_{\min}+(\beta_{\max}-\beta_{\min})\odot U] \]
    \emphok{Điểm mấu chốt:} nhiễu lấy $\mu$ và $\sigma$ của chính tập dữ liệu,
    nên vector nhiễu vẫn nằm trong vùng hợp lý của không gian đặc trưng.
  \end{methodblock}
\end{column}
\begin{column}{0.47\textwidth}
  \begin{contextblock}{Quyết định có tiêm hay không}
    \[
      \mathbf{V}_n=
      \begin{cases}
        (1-\theta_r)\mathbf{V}_o+\mathbf{N}_m & \text{nếu } \rho>\theta_n\\
        \mathbf{V}_o & \text{ngược lại}
      \end{cases}
    \]
    với $\rho\sim U[0,1]$ cho từng thực thể.
  \end{contextblock}

  \begin{riskblock}{Đối chiếu với mã nguồn}
    \texttt{\scriptsize self.mask\_ratio\ =\ nn.Parameter(0.7)}\\
    \texttt{\scriptsize self.noise\_ratio = nn.Parameter(0.2)}
    \\[3pt]
    Bài báo mô tả $\theta_n$ là \emph{``adaptively tuned during training''}.
    \emphwarn{Nhóm đã kiểm chứng lại điểm này} --- xem phần thực nghiệm.
  \end{riskblock}
\end{column}
\end{columns}
\end{frame}

\begin{frame}{Đóng góp 2 --- Attention Penalty: ba nhịp}
\begin{columns}[T]
\begin{column}{0.32\textwidth}
  \begin{riskblock}{Nhịp 1: phát hiện}
    Tích luỹ hậu tố của điểm chú ý:
    \[ P[\dots,i]=\textstyle\prod_{j\ge i} A[\dots,j] \]
    Vượt ngưỡng $\phi=10$ thì bị đánh dấu.
    \\[3pt]
    \emphwarn{Vì sao dùng tích?} Phép nhân khuếch đại: chỉ token cao
    \emph{dai dẳng} mới bị bắt, cao \emph{nhất thời} thì không.
  \end{riskblock}
\end{column}
\begin{column}{0.32\textwidth}
  \begin{riskblock}{Nhịp 2: phạt}
    \[ A^p_j = A^p_i-\theta_p(A^p_i),\ \text{nếu } A^p_i>\phi \]
    Trong cài đặt: \texttt{clamp} giới hạn biên độ rồi nhân
    \texttt{penalty\_weights} $=0.001$ để không phạt quá tay.
  \end{riskblock}
\end{column}
\begin{column}{0.32\textwidth}
  \begin{methodblock}{Nhịp 3: tái phân bổ}
    \[ A^n_j = A^n_i+\theta^n_j(A^p_i-\phi) \]
    với $n\sim(p-k,\,p+k)$, $k=1$.
    \\[3pt]
    Phần chú ý bị tịch thu \emphok{không biến mất} mà chia lại cho token lân cận
    --- đây là lý do bài báo gọi là \emph{redistribution}.
  \end{methodblock}
\end{column}
\end{columns}

\vspace{4pt}
\begin{objectiveblock}{Vì sao cần cả tăng lẫn giảm --- bằng chứng từ ablation}
  Bảng 4 bỏ riêng từng nhánh: bỏ nhánh tăng cho MRR $0.365$, bỏ nhánh giảm cho
  $0.369$, đủ cả hai cho $\mathbf{0.374}$. Nghĩa là \emphobj{chỉ phạt thôi chưa
  đủ} --- phải trả chú ý về đúng chỗ thì mô hình mới hưởng lợi.
\end{objectiveblock}
\end{frame}

\begin{frame}{Hợp nhất mô thái và chấm điểm}
\begin{columns}[T]
\begin{column}{0.52\textwidth}
  \begin{methodblock}{Modality interaction}
    \[ \theta_{ab}=\frac{\exp(Q_a^{\!\top}K_b/\sqrt{d})}
                        {\exp(Q_a^{\!\top}K_a/\sqrt{d})+\exp(Q_a^{\!\top}K_b/\sqrt{d})} \]
    \[ \mathbf{V}^i_a=\theta_{ab}\mathbf{V}_a \]
    Khác multi-head attention thường: công thức này \emphok{cân giữa tự tương
    quan và tương quan chéo}. Nếu ảnh của thực thể mờ, $\theta$ tự nghiêng về
    văn bản.
  \end{methodblock}

  \begin{contextblock}{Hàm mất mát}
    \[ \mathcal{L}=\frac{1}{|\mathcal{T}|}\sum\Big[-\log\sigma(\lambda-S)
       -\sum_i v_i\,\sigma(S'-\lambda)\Big] \]
    $v_i$ là trọng số tự đối kháng, kế thừa từ RotatE.
  \end{contextblock}
\end{column}
\begin{column}{0.45\textwidth}
  \begin{contextblock}{Sáu phương án hợp nhất trong mã}
    {\scriptsize
    \begin{tabular}{@{}ll@{}}
      \toprule
      Phương án & Dùng cho\\
      \midrule
      \texttt{Mformer\_hd\_graph} & \textbf{DB15K}\\
      \texttt{Mformer\_hd\_mean}  & \textbf{MKG-W}\\
      \texttt{Mformer\_weight}    & ---\\
      \texttt{atten\_weight}      & ---\\
      \texttt{learnable\_weight}  & ---\\
      \texttt{cat\_mlp}           & ---\\
      \bottomrule
    \end{tabular}}
  \end{contextblock}

  \begin{riskblock}{Chi tiết đáng chú ý}
    Hai tập dữ liệu dùng \emphwarn{hai chiến lược khác nhau} cho cả nhiễu lẫn
    hợp nhất. Vừa là thích ứng hợp lý, vừa làm dấy lên nghi vấn về mức độ điều
    chỉnh riêng cho từng tập.
  \end{riskblock}
\end{column}
\end{columns}
\end{frame}

% ==============================================================================
\section{Ví dụ chạy qua mô hình}
% ==============================================================================

\begin{frame}{Một bộ ba đi xuyên mô hình}
\begin{contextblock}{Đầu vào}
  \texttt{\scriptsize ( National Society of Film Critics Award for Best Supporting
  Actor,\ \ award,\ \ \textcolor{BadRed}{?} )}
  \\[2pt]
  Kèm ảnh logo giải thưởng và mô tả văn bản từ DBpedia. Đáp án đúng:
  \emphkey{Dennis Hopper}.
\end{contextblock}

\vspace{2pt}
\begin{columns}[T]
\begin{column}{0.32\textwidth}
  {\scriptsize\textbf{Bước 1--2: tra cứu, tiêm nhiễu}}\\[2pt]
  {\tiny\ttfamily
   h\_id = 3417\\
   $\es$ = ent\_emb[3417], 256-d\\
   $\ev$ = BEiT, 4096-d\\
   $\et$ = BERT, 768-d\\
   \textcolor{OrangeAc}{$\rho$ = 0.31 > $\theta_n$ = 0.2}\\
   nên tiêm nhiễu vào $\ev$}
\end{column}
\begin{column}{0.32\textwidth}
  {\scriptsize\textbf{Bước 3: chiếu, self-attention}}\\[2pt]
  {\tiny\ttfamily
   img\_proj: 4096 $\to$ 256\\
   text\_proj: 768 $\to$ 256\\
   stack thành [B, 3, 256]\\
   A = QK$^{\top}$/$\sqrt{d}$\\
   \textcolor{BadRed}{A["for"] = 0.66} (over-trust)}
\end{column}
\begin{column}{0.32\textwidth}
  {\scriptsize\textbf{Bước 4: phạt, tái phân bổ}}\\[2pt]
  {\tiny\ttfamily
   P = suffix\_prod(A)\\
   P["for"] = 12.4 > $\phi$ = 10\\
   nên bị phạt\\
   \textcolor{GoodGreen}{A["for"]: 0.66 $\to$ 0.50}\\
   phần dư chia cho token cạnh}
\end{column}
\end{columns}

\vspace{6pt}
\begin{columns}[T]
\begin{column}{0.48\textwidth}
  \begin{riskblock}{Không có attention penalty}
    1.\ Kurtwood Smith\\
    2.\ \emphwarn{Dennis Hopper} --- đúng nhưng xếp nhì\\
    3.\ James Marsden
  \end{riskblock}
\end{column}
\begin{column}{0.48\textwidth}
  \begin{methodblock}{Có attention penalty}
    1.\ \emphok{Dennis Hopper} --- lên hạng nhất\\
    2.\ Kurtwood Smith\\
    3.\ James Marsden
  \end{methodblock}
\end{column}
\end{columns}
\end{frame}

% ==============================================================================
\section{Thực nghiệm trong bài báo}
% ==============================================================================

\begin{frame}{So sánh với các mô hình nền}
\begin{columns}[T]
\begin{column}{0.53\textwidth}
  {\scriptsize
  \begin{tabular}{@{}lcccc@{}}
    \toprule
    \multirow{2}{*}{Mô hình} & \multicolumn{2}{c}{DB15K} & \multicolumn{2}{c}{MKG-W}\\
    \cmidrule(lr){2-3}\cmidrule(lr){4-5}
     & MRR & H@1 & MRR & H@1\\
    \midrule
    TransE \tiny(NeurIPS'13) & 0.249 & 0.128 & 0.292 & 0.211\\
    RotatE \tiny(ICLR'19)    & 0.293 & 0.179 & 0.337 & 0.268\\
    PairRE \tiny(ACL'21)     & 0.311 & 0.216 & 0.344 & 0.282\\
    IMF \tiny(WWW'23)        & 0.323 & 0.242 & 0.345 & 0.288\\
    AdaMF \tiny(COLING'24)   & 0.325 & 0.213 & 0.343 & 0.272\\
    SnAg \tiny('24)          & 0.342 & 0.250 & 0.359 & 0.284\\
    NativE \tiny(SIGIR'24)   & 0.362 & 0.273 & 0.366 & 0.296\\
    \midrule
    \rowcolor{BlueBg}
    \textbf{APKGC} & \textbf{0.364} & \textbf{0.282} & \textbf{0.374} & \textbf{0.306}\\
    \bottomrule
  \end{tabular}}
\end{column}
\begin{column}{0.44\textwidth}
  \begin{contextblock}{Điều bài báo nhấn mạnh}
    Hits@1 tăng \emphkey{3.3\%} trên DB15K ($0.273\to0.282$) và \emphkey{3.4\%}
    trên MKG-W ($0.296\to0.306$) so với NativE.
  \end{contextblock}

  \begin{riskblock}{Điều cần nói thẳng}
    \begin{itemize}\itemsep2pt
      \item Đây là mức tăng \emphwarn{tương đối}; chênh lệch tuyệt đối chỉ khoảng
            \emphwarn{0.9 điểm phần trăm}.
      \item Ở Hits@10 trên DB15K, APKGC \emphwarn{thua} NativE ($0.527$ so với
            $0.529$) và bài báo không thảo luận.
      \item Bảng \emphwarn{không báo cáo độ lệch chuẩn}. Nhóm đã tự đo --- xem
            phần sau.
    \end{itemize}
  \end{riskblock}
\end{column}
\end{columns}
\end{frame}

\begin{frame}{Ablation: bỏ thành phần nào thì hỏng đến đâu}
\begin{columns}[T]
\begin{column}{0.48\textwidth}
  {\scriptsize
  \textbf{Theo mô thái (DB15K)}\\[3pt]
  \begin{tabular}{@{}lccc@{}}
    \toprule
    Cấu hình & MRR & H@1 & H@10\\
    \midrule
    Chỉ cấu trúc       & 0.211 & 0.125 & 0.345\\
    Chỉ ảnh            & 0.336 & 0.261 & 0.469\\
    Chỉ văn bản        & 0.348 & 0.273 & 0.505\\
    Cấu trúc + ảnh     & 0.340 & 0.260 & 0.480\\
    Văn bản + cấu trúc & 0.351 & 0.279 & 0.509\\
    Ảnh + văn bản      & 0.357 & 0.274 & 0.521\\
    \midrule
    \rowcolor{BlueBg}
    \textbf{Đủ ba}     & \textbf{0.364} & \textbf{0.282} & \textbf{0.527}\\
    \bottomrule
  \end{tabular}}
  \\[4pt]
  \muted{\scriptsize Chỉ cấu trúc kém xa phần còn lại. Văn bản đóng góp nhiều hơn
  ảnh: bỏ văn bản mất 8.9\% H@10, bỏ ảnh chỉ mất 3.4\%.}
\end{column}
\begin{column}{0.48\textwidth}
  {\scriptsize
  \textbf{Theo thành phần (MKG-W)}\\[3pt]
  \begin{tabular}{@{}lccc@{}}
    \toprule
    Cấu hình & MRR & H@1 & H@10\\
    \midrule
    Bỏ lấy mẫu nhiễu     & 0.358 & 0.289 & 0.486\\
    Bỏ attention penalty & 0.360 & 0.286 & 0.492\\
    Bỏ nhánh tăng        & 0.365 & 0.297 & 0.491\\
    Bỏ nhánh giảm        & 0.369 & 0.302 & 0.495\\
    \midrule
    \rowcolor{BlueBg}
    \textbf{Đầy đủ}      & \textbf{0.374} & \textbf{0.306} & \textbf{0.501}\\
    \bottomrule
  \end{tabular}}
  \\[4pt]
  \begin{objectiveblock}{Quan sát thú vị}
    \emphobj{Lấy mẫu nhiễu đóng góp nhiều hơn attention penalty} (mất 0.016 so
    với 0.014 MRR) --- dù tiêu đề bài báo đặt attention penalty ở vị trí nổi bật
    hơn.
  \end{objectiveblock}
\end{column}
\end{columns}
\end{frame}

\begin{frame}{Phân tích siêu tham số: phạt quá tay cũng hỏng}
\begin{columns}[c]
\begin{column}{0.46\textwidth}
  \centering
  {\scriptsize DB15K --- MRR theo số cửa sổ $k$ và tỉ lệ phạt $\theta_p$}\\[4pt]
  \begin{tikzpicture}[scale=0.9, transform shape,
    lb/.style={font=\tiny, text=MutedTxt}]
    \node[lb] at (-0.5,2.95) {$k$=4};
    \node[lb] at (-0.5,2.25) {$k$=3};
    \node[lb] at (-0.5,1.55) {$k$=2};
    \node[lb] at (-0.5,0.85) {$k$=1};
    \node[lb] at (0.60,0.20) {1};
    \node[lb] at (1.85,0.20) {0.001};
    \node[lb] at (3.10,0.20) {0.01};
    \node[lb] at (4.35,0.20) {0.1};
    \node[lb] at (2.4,-0.20) {tỉ lệ phạt $\theta_p$};

    \fill[structblue!55] (0.0,2.65) rectangle ++(1.2,0.6);
    \fill[structblue!54] (1.25,2.65) rectangle ++(1.2,0.6);
    \fill[structblue!38] (2.50,2.65) rectangle ++(1.2,0.6);
    \fill[structblue!33] (3.75,2.65) rectangle ++(1.2,0.6);
    \node[font=\tiny] at (0.60,2.95) {.345};
    \node[font=\tiny] at (1.85,2.95) {.343};
    \node[font=\tiny] at (3.10,2.95) {.329};
    \node[font=\tiny] at (4.35,2.95) {.324};

    \fill[structblue!58] (0.0,1.95) rectangle ++(1.2,0.6);
    \fill[structblue!62] (1.25,1.95) rectangle ++(1.2,0.6);
    \fill[structblue!52] (2.50,1.95) rectangle ++(1.2,0.6);
    \fill[structblue!42] (3.75,1.95) rectangle ++(1.2,0.6);
    \node[font=\tiny] at (0.60,2.25) {.348};
    \node[font=\tiny] at (1.85,2.25) {.352};
    \node[font=\tiny] at (3.10,2.25) {.342};
    \node[font=\tiny] at (4.35,2.25) {.332};

    \fill[structblue!60] (0.0,1.25) rectangle ++(1.2,0.6);
    \fill[structblue!66] (1.25,1.25) rectangle ++(1.2,0.6);
    \fill[structblue!95] (2.50,1.25) rectangle ++(1.2,0.6);
    \fill[structblue!43] (3.75,1.25) rectangle ++(1.2,0.6);
    \node[font=\tiny] at (0.60,1.55) {.350};
    \node[font=\tiny] at (1.85,1.55) {.356};
    \node[font=\tiny\bfseries, text=white] at (3.10,1.55) {.364};
    \node[font=\tiny] at (4.35,1.55) {.330};

    \fill[structblue!44] (0.0,0.55) rectangle ++(1.2,0.6);
    \fill[structblue!62] (1.25,0.55) rectangle ++(1.2,0.6);
    \fill[structblue!66] (2.50,0.55) rectangle ++(1.2,0.6);
    \fill[structblue!55] (3.75,0.55) rectangle ++(1.2,0.6);
    \node[font=\tiny] at (0.60,0.85) {.330};
    \node[font=\tiny] at (1.85,0.85) {.349};
    \node[font=\tiny] at (3.10,0.85) {.353};
    \node[font=\tiny] at (4.35,0.85) {.344};
  \end{tikzpicture}
\end{column}
\begin{column}{0.5\textwidth}
  \begin{objectiveblock}{Đường cong hình chữ U ngược}
    Hiệu năng đạt đỉnh ở \emphobj{vùng giữa} ($k=2$, $\theta_p=0.01$, MRR
    $0.364$) rồi suy giảm về cả hai phía.
  \end{objectiveblock}

  \begin{riskblock}{Hai đầu đều xấu}
    \emphwarn{Phạt quá nhẹ:} không đủ kéo chú ý khỏi token vô nghĩa.\\[3pt]
    \emphwarn{Phạt quá nặng:} phá vỡ cả những phân bổ chú ý đúng đắn. Tại
    $\theta_p=0.1$, MRR rơi xuống $0.324$ --- thấp hơn cả khi không dùng penalty.
  \end{riskblock}

  \begin{contextblock}{Ý nghĩa thực tiễn}
    Đây là hạn chế về khả năng triển khai: mỗi tập dữ liệu mới đều đòi hỏi dò
    lại lưới siêu tham số, và bài báo không đề xuất cách chọn tự động.
  \end{contextblock}
\end{column}
\end{columns}
\end{frame}

% ==============================================================================
\section{Thực nghiệm của nhóm}
% ==============================================================================

\begin{frame}{Chúng tôi chạy lại --- số liệu tự đo}
\begin{columns}[T]
\begin{column}{0.5\textwidth}
  {\scriptsize
  \textbf{MKG-W --- Kaggle, GPU Tesla P100}\\[4pt]
  \begin{tabular}{@{}lcccc@{}}
    \toprule
    Seed & MRR & H@1 & H@3 & H@10\\
    \midrule
    \multicolumn{5}{@{}l}{\itshape Đợt 1 --- 8000 epoch, đánh giá mỗi 100}\\
    7    & 0.3616 & 0.2942 & 0.3918 & 0.4853\\
    42   & 0.3641 & 0.2956 & 0.3954 & 0.4840\\
    \multicolumn{5}{@{}l}{\itshape Đợt 2 --- 2000 epoch, đánh giá mỗi 250}\\
    1234 & 0.3610 & 0.2987 & 0.3864 & 0.4759\\
    2024 & 0.3611 & 0.2943 & 0.3907 & 0.4819\\
    3407 & 0.3595 & 0.2925 & 0.3917 & 0.4799\\
    \midrule
    \rowcolor{GoodBg}
    \textbf{TB đợt 2} & \textbf{0.3605} & 0.2952 & 0.3896 & 0.4792\\
    Bài báo & 0.3740 & 0.3060 & 0.4040 & 0.5010\\
    \bottomrule
  \end{tabular}}
  \\[5pt]
  \muted{\scriptsize Hai đợt \textbf{khác điều kiện} nên báo cáo tách biệt. Ước
  lượng độ lệch chuẩn lấy từ đợt 2 --- ba lần chạy duy nhất cùng điều kiện.}
  \\[4pt]
  \emphok{\scriptsize Tái lập: chạy lại seed 42 sau sáu ngày, phiên độc lập ---
  bốn chỉ số \textbf{trùng đến chữ số thứ năm}. Quy trình có tính tất định.}
\end{column}
\begin{column}{0.47\textwidth}
  \begin{methodblock}{Mô hình ổn định giữa các seed}
    Độ lệch chuẩn MRR (đợt 2, $n{=}3$) $=\mathbf{0.0009}$.
    \\[3pt]
    {\scriptsize
    \begin{tabular}{@{}lcc@{}}
      Khoảng cách trong bài báo & giá trị & so với std\\
      \midrule
      vs NativE (DB15K) & 0.002 & $2.2\times$\\
      vs NativE (MKG-W) & 0.008 & $8.9\times$\\
      attention penalty & 0.014 & $15.6\times$\\
      noise sampling    & 0.016 & $17.8\times$\\
    \end{tabular}}
    \\[3pt]
    \emphok{Mọi khoảng cách bài báo báo cáo đều vượt nhiễu seed. Giả thuyết ban
    đầu của nhóm bị chính số liệu bác bỏ.}
  \end{methodblock}

  \begin{riskblock}{Lệch hệ thống $-0.0135$ MRR}
    Đều ở cả bốn chỉ số, \emphwarn{cả năm lần chạy} --- gấp \emphwarn{mười lăm
    lần} độ lệch chuẩn, nên không phải dao động ngẫu nhiên.
    \\[3pt]
    Nghi vấn: bài báo ghi batch 2048 / neg 64, script chạy 1024 / 32.
    \emphwarn{Nhóm đã chạy thí nghiệm tách nguyên nhân} --- xem slide sau.
    \\[3pt]
    {\scriptsize Không do dao động khi chạy: dao động đó \textbf{bằng không}.}
  \end{riskblock}
\end{column}
\end{columns}
\end{frame}

\begin{frame}{Tách nguyên nhân: chạy đúng cấu hình bài báo}
\begin{columns}[T]
\begin{column}{0.52\textwidth}
  {\scriptsize
  \textbf{Chạy lại với batch 2048 / neg 64} (đúng tr.\ 15010), seed 42\\[2pt]
  {\tiny Cùng điều kiện đợt 2 nên đối chứng là đợt 2, không phải trung bình gộp.}\\[4pt]
  \begin{tabular}{@{}lccc@{}}
    \toprule
    Chỉ số & Đợt 2 & Bài báo & Công bố\\
           & \scriptsize 1024/32 & \scriptsize 2048/64 & \\
    \midrule
    MRR    & 0.3605 & \textbf{0.3691} & 0.3740\\
    H@1    & 0.2952 & \textbf{0.3028} & 0.3060\\
    H@3    & 0.3896 & \textbf{0.3997} & 0.4040\\
    H@10   & 0.4792 & 0.4874 & 0.5010\\
    \bottomrule
  \end{tabular}
  \\[8pt]
  \textbf{Khoảng cách tới số công bố}\\[4pt]
  \begin{tabular}{@{}lccc@{}}
    \toprule
    Chỉ số & Đợt 2 & Bài báo & Thu hẹp\\
    \midrule
    MRR  & $-0.0135$ & $-0.0049$ & \textbf{64\%}\\
    H@1  & $-0.0108$ & $-0.0032$ & \textbf{70\%}\\
    H@3  & $-0.0144$ & $-0.0043$ & \textbf{70\%}\\
    H@10 & $-0.0218$ & $-0.0136$ & 37\%\\
    \bottomrule
  \end{tabular}}
\end{column}
\begin{column}{0.45\textwidth}
  \begin{methodblock}{Giả thuyết được xác nhận}
    Chạy đúng cấu hình ghi trong bài báo \emphok{thu hẹp 64\%} khoảng lệch MRR,
    và \emphok{70\%} ở Hits@1 lẫn Hits@3. Riêng cấu hình đóng góp
    \emphok{$+0.0086$ MRR} --- gấp gần mười lần độ lệch chuẩn.
    \\[3pt]
    Phát biểu kiểm chứng được: \textit{tệp \texttt{run.sh} gọi script đặt batch
    1024 / neg 32, khác con số 2048 / 64 ghi trong bài. Người làm theo hướng dẫn
    trong repo sẽ không tái tạo được bảng kết quả.}
  \end{methodblock}

  \begin{riskblock}{Phần chưa giải thích: 0.0049 MRR}
    Vẫn gấp năm lần độ lệch chuẩn. Ba nguyên nhân còn lại:
    \begin{itemize}\itemsep1pt
      \item Phần cứng: P100 so với RTX 4090
      \item torch 2.4.1 thay vì bản tác giả dùng
      \item Chỉ chạy 2000 epoch, đánh giá mỗi 250 (script gốc: 8000 và 100)
    \end{itemize}
    \emphwarn{Chỉ một seed} nên chưa có độ lệch chuẩn riêng.
  \end{riskblock}
\end{column}
\end{columns}
\end{frame}

\begin{frame}{Sáu rào cản trước khi chạy được dòng đầu tiên}
\begin{columns}[T]
\begin{column}{0.5\textwidth}
  \begin{riskblock}{Trong mã nguồn công bố}
    \begin{itemize}\itemsep2pt
      \item \emphwarn{Thiếu hẳn mã C++.} \texttt{make.sh} cần
            \texttt{mmkgc/base/} nhưng thư mục không được commit;
            \texttt{Base.so} sẵn có là ELF x86-64.
      \item \emphwarn{17 lời gọi \texttt{.cuda()} cứng} --- không có GPU NVIDIA
            là hỏng.
      \item \texttt{requirement.txt} thiếu \texttt{scikit-learn}, \texttt{scipy}.
      \item Ghi kết quả ra \texttt{../../data/} --- ngoài thư mục dự án.
      \item Sập nếu chạy dưới 100 epoch: \texttt{best\_model\_wts} truy cập vô
            điều kiện.
      \item \texttt{torch>=1.6.0} không có cận trên --- bản mới bỏ hỗ trợ sm\_60.
    \end{itemize}
  \end{riskblock}
\end{column}
\begin{column}{0.47\textwidth}
  \begin{methodblock}{Cách nhóm khắc phục}
    \texttt{mmkgc/} hoá ra là bản sửa đổi của \emphok{OpenKE} (THUNLP). Lấy
    \texttt{base/} từ repo gốc, đối chiếu chữ ký \texttt{sampling()} --- khớp
    đúng 11 tham số --- rồi biên dịch lại. Chạy được trên cả ARM64 lẫn x86-64.
  \end{methodblock}

  \begin{objectiveblock}{Chi phí tính toán --- nút thắt bất ngờ}
    {\scriptsize
    \begin{tabular}{@{}lr@{}}
      \toprule
      Khâu & Thời gian\\
      \midrule
      Huấn luyện & 8.4 giây / epoch\\
      \textbf{Đánh giá} & \textbf{21 phút / lần}\\
      \bottomrule
    \end{tabular}}
    \\[4pt]
    \texttt{run\_link\_prediction} xếp hạng \emphobj{tuần tự} 4.274 bộ ba, hai
    chiều, mỗi lần trên 15.000 thực thể, không gộp lô. Cấu hình gốc cần khoảng
    \emphobj{28 giờ chỉ riêng đánh giá}.
  \end{objectiveblock}
\end{column}
\end{columns}
\end{frame}

% ==============================================================================
\section{Đánh giá và hướng cải tiến}
% ==============================================================================

\begin{frame}{Đọc bài báo bằng con mắt phản biện}
\begin{columns}[T]
\begin{column}{0.48\textwidth}
  \begin{methodblock}{Điểm mạnh}
    \begin{itemize}\itemsep2pt
      \item \emphok{Lật ngược một giả định phổ biến.} Nhiễu từ kẻ thù thành công
            cụ --- đóng góp về tư duy, không chỉ kỹ thuật.
      \item \emphok{Chẩn đoán đúng vấn đề.} Attention over-trust được ghi nhận
            độc lập trong nghiên cứu khác.
      \item \emphok{Tái phân bổ chứ không chỉ triệt tiêu.} Ablation chứng minh
            cả hai nhánh đều đóng góp.
      \item \emphok{Ổn định giữa các seed} --- nhóm tự kiểm chứng, độ lệch
            chuẩn chỉ $0.0009$.
    \end{itemize}
  \end{methodblock}
\end{column}
\begin{column}{0.48\textwidth}
  \begin{riskblock}{Hạn chế}
    \begin{itemize}\itemsep2pt
      \item Mức cải thiện dưới một điểm phần trăm tuyệt đối, \emphwarn{không có
            độ lệch chuẩn} trong bài.
      \item \emphwarn{Thua ở Hits@10} trên DB15K nhưng không thảo luận.
      \item Ngưỡng cứng thiếu biện luận: $\phi{=}10$, \texttt{max}${=}15$,
            \texttt{weights}${=}0.001$.
      \item Cấu hình khác nhau theo từng tập dữ liệu.
      \item \emphwarn{Bài báo và mã nguồn không khớp}: batch 2048 / neg 64 so với
            1024 / 32 trong script.
      \item Tích luỹ hậu tố dùng vòng lặp Python --- nhóm đã thử vector hoá,
            hoá ra \emphwarn{không có lợi} ở quy mô hiện tại.
    \end{itemize}
  \end{riskblock}
\end{column}
\end{columns}
\end{frame}

\begin{frame}{Sáu hướng cải tiến, xếp theo công sức}
{\scriptsize
\begin{tabular}{@{}p{2.2cm}p{4.0cm}p{6.4cm}@{}}
  \toprule
  Mức độ & Đề xuất & Nội dung\\
  \midrule
  \textcolor{MutedTxt}{\textbf{Đã thử}} \muted{xong}
    & Vector hoá tích luỹ hậu tố
    & Đã kiểm chứng: kết quả trùng khớp (sai lệch $3{,}8\!\times\!10^{-6}$) nhưng
      \emphwarn{chậm hơn 1,4 lần} ở cấu hình hiện tại vì chiều cuối chỉ bằng 3.
      Chỉ có lợi khi chuỗi dài từ 8 trở lên.\\
  \addlinespace[2pt]
  \textcolor{GoodGreen}{\textbf{Dễ}} \muted{một ngày}
    & Ngưỡng thích ứng
    & Đặt $\phi$ theo phân vị của phân phối chú ý trong từng lô thay vì cố định
      bằng 10. Hình 3 cho thấy hiệu năng rất nhạy với ngưỡng.\\
  \addlinespace[2pt]
  \textcolor{OrangeAc}{\textbf{Trung bình}} \muted{vài ngày}
    & Lịch trình nhiễu
    & Giảm dần cường độ nhiễu theo tiến trình huấn luyện --- mượn ý tưởng từ
      curriculum learning và lịch trình nhiễu của mô hình khuếch tán.\\
  \addlinespace[2pt]
  \textcolor{OrangeAc}{\textbf{Trung bình}} \muted{một tuần}
    & Chạy đúng cấu hình bài báo
    & Batch 2048, neg 64 --- để tách riêng nguyên nhân của khoảng lệch $0.011$
      mà nhóm quan sát được.\\
  \addlinespace[2pt]
  \textcolor{BadRed}{\textbf{Khó}} \muted{một tháng}
    & Kiểm chứng trên tập thứ ba
    & Bài báo chỉ thử hai tập. Chạy trên MKG-Y trả lời câu hỏi lớn nhất: mô hình
      có tổng quát hoá được không?\\
  \addlinespace[2pt]
  \textcolor{objectivepurple}{\textbf{Mở rộng}}
    & Thêm mô thái thứ tư
    & Chính tác giả gợi ý đưa video và âm thanh vào. Cơ chế tương quan chéo
      $\theta_{ab}$ về nguyên tắc mở rộng được cho $n$ mô thái.\\
  \bottomrule
\end{tabular}}
\end{frame}

\begin{frame}{Ba điều mang về}
\begin{columns}[T]
\begin{column}{0.32\textwidth}
  \begin{contextblock}{Một}
    Nhiễu có thể là \emphkey{bạn}. Tiêm nhiễu có kiểm soát trong lúc huấn luyện
    làm mô hình bền vững hơn là cố lọc sạch dữ liệu --- cùng tinh thần với
    dropout và data augmentation, nay áp dụng ở mức mô thái.
  \end{contextblock}
\end{column}
\begin{column}{0.32\textwidth}
  \begin{riskblock}{Hai}
    Attention \emphwarn{không tự biết} mình đang nhìn nhầm chỗ. Can thiệp trực
    tiếp vào ma trận chú ý --- phạt rồi trả lại --- là hướng chỉnh sửa hiếm gặp
    trong các công trình MMKGC.
  \end{riskblock}
\end{column}
\begin{column}{0.32\textwidth}
  \begin{methodblock}{Ba}
    \emphok{Đọc mã nguồn cho biết những gì đọc bài báo không cho biết.} Nhóm
    phát hiện sáu rào cản kỹ thuật, và xác nhận bằng thực nghiệm rằng
    \textbf{64\% khoảng lệch} đến từ việc script công bố không khớp cấu hình
    ghi trong bài.
  \end{methodblock}
\end{column}
\end{columns}

\vspace{8pt}
\begin{center}
  \muted{\scriptsize Tài liệu kèm theo: báo cáo chi tiết $\cdot$ kết quả chạy lại
  $\cdot$ rà soát tài liệu $\cdot$ notebook tái lập}
\end{center}
\end{frame}

\end{document}
